\documentclass[10pt,a4paper]{amsart}
\usepackage[margin=19mm]{geometry}
\usepackage{amsmath,amssymb,mathtools}
\usepackage[scr=rsfs,cal=euler]{mathalfa}
\usepackage{xfrac}
\usepackage[unicode,hidelinks,bookmarks=false]{hyperref}
\newcommand{\ie}{i.e.\ }
\newcommand{\cf}{cf.\ }
\newcommand{\resp}{respectively\ }
\newcommand{\Subsection}{\subsection}
\providecommand{\nobreakdash}{\nobreak\hbox{-}}
\setlength{\emergencystretch}{2em}
\multlinegap0pt
\usepackage{bm}
\usepackage{amssymb}
\usepackage{dsfont}
\usepackage{mathtools}
\newtheorem{corollary}{Corollary}
\newtheorem{theorem}{Theorem}
\newcommand{\A}{\mathbb A}
\DeclareMathOperator{\Gal}{Gal}
\newcommand{\I}{\mathbb{I}}
\DeclareMathOperator{\Pic}{Pic}
\newcommand{\Q}{\mathbb Q}
\newcommand{\Z}{\mathbb Z}
\newcommand{\cor}{\mathrm{cor}}
\newcommand{\ord}{\mathrm{ord}}
\newcommand{\res}{\mathrm{res}}
\title{Quaternionic families of Heegner points and $p$-adic $L$-functions}
\author{Matteo Longo, Paola Magrone, Eris Rocha Walchek}
\date{Source: arXiv:2510.04306v3 (CC BY 4.0)}
\begin{document}
\maketitle

\noindent\emph{Source and licence.} Quaternionic families of Heegner points and $p$-adic $L$-functions. Version arXiv:2510.04306v3 (CC BY 4.0), distributed by arXiv under the Creative Commons Attribution 4.0 International licence, http://creativecommons.org/licenses/by/4.0/, as recorded in the arXiv licence metadata for this exact version and shown on the abstract page https://arxiv.org/abs/2510.04306v3. Full licence notice: \texttt{licenses/arxiv-2510.04306-CC-BY-4.0.txt}.\par\smallskip

\noindent\textit{Editorial scope.} Counted unit: the article's theorem thm13.5 (source0.tex lines 1395--1401). The article's complete original proof of it is delivered with it (source0.tex lines 1403--1473, one \texttt{proof} environment of 71 lines). No mathematics is written, repaired or re-derived: the statement and the proof are the article's own lines, in the article's own notation, and no retained line is substituted or re-flowed.  Retained uncounted as the source-local context the proof needs: lines 747--751; lines 1125--1126; lines 1276--1278; lines 1379--1380; lines 1388--1390.  Nothing was omitted from any retained span, so the package carries no omission record.  No retained line contains a citation, so the wrapper carries no bibliography and no bibliography entry.  No retained line contains a figure environment, an image-inclusion command or drawing code, so no image is delivered and the package declares no asset dependency.  Editorial formatting only, in addition: an AMS \texttt{amsart} preamble that restates the article's own declarations and macros used by the retained lines, namely the environments \texttt{corollary}, \texttt{theorem} on separate counters, together with the standard packages those lines need.  The retained environments print in this document as Corollary 1, Theorem 1 in order of appearance, whereas the article prints the same items with the numbering of its own full document.  The complete native source of the article as distributed in the arXiv e-print for this exact version is bundled unchanged as \texttt{sources/186431/source0.tex}.
\begin{equation}\label{prop12.2}
\begin{split}
\mathrm{sp}_{\nu,\phi^{-1}}\left(\mathcal{L}_{\omega_{\I}}^{\Gamma_\infty}(\mathfrak{Y})\right)&=\frac{(-1)^{-w-1}}{(-w-1)!}\cdot \mathcal{E}(\phi,\nu)\cdot
({\omega}_{\mathcal{F}_\nu}\otimes\phi)\left(\log(\mathrm{sp}_{\nu,\phi^{-1}}(\mathfrak{Y}))\right)
.\end{split}\end{equation}

\begin{equation}\label{Col-dminus1}
d^{-1}\mathcal{F}^{[p]}_{x_\infty}=\Pi(\phi_{\mathrm{DT}}^*)F_\infty.\end{equation} 

\begin{corollary}\label{CoroCole} Let $\Delta=(P)-(x_\infty)$ and $F_\infty^*$ the Coleman primitive of $\omega_{\mathcal{F}^*}$ on $\mathcal{W}_\infty$ which vanishes at $\infty$. Assume that $m>1$. Then 
$\delta_m(\Delta)(\omega_{\mathcal{F}^*})=F_{\omega_{\mathcal{F}^*}}(P)$. 
\end{corollary} 

\begin{equation}\label{eq9.14}{J}_m(L)^\ord\otimes_\Z\mathcal{O}\stackrel{\sim}{\longrightarrow} 
\Pic({X}_m/L)^\ord\otimes_\Z\mathcal{O}.\end{equation} 

\begin{equation}\label{eq13.2}
\varrho_m(Q_{cp^n,m})=\left(\frac{\nu(\mathbf{a}_p)}{p}\right)^m\cdot \mathrm{sp}_\nu(\mathfrak{X}_{cp^n}). 
\end{equation} 

\begin{theorem}\label{thm13.5} Let $\nu$ be an arithmetic morphism with signature $(2,\psi)$, where $\mathrm{cond}(\psi)=p^m$ and $m\geq 2$, and $\phi\colon K^\times\backslash\A_K^\times\rightarrow\overline\Q^\times$ be of infinity type $(1,-1)$ and $\mathrm{cond}(\phi)=p^n$ with $n\geq m$. 
Then 
\[\mathscr{L}_{\I,\boldsymbol{\xi}}^{\mathrm{alg}}(\nu,\hat\phi^{-1})=\frac{\epsilon(\phi)}
{{\xi}_{\nu,\mathfrak{p}}(p^n)\cdot p^n}\cdot\sum_{\mathfrak{a}\in\Pic(\mathcal{O}_{cp^n})}
(\hat\xi_\nu^{-1}\hat\chi_\nu\hat\phi)(a)d^{-1}\mathcal{F}_{\nu , x_\infty}^{[p]}\left(x_{cp^n,m}(\mathfrak{a}^{-1})\right)
.\] 
\end{theorem}

\begin{proof} We first relate $\mathscr{L}_{\I,\boldsymbol{\xi}}^{\mathrm{alg}}(\nu,\hat\phi^{-1})$ to the Coleman primitive. 
Since $\hat\phi\colon \Gamma_\infty\rightarrow\overline\Q_p$ 
has Hodge--Tate weight $w=1$ (so $\hat{\phi}^{-1}$ has Hodge--Tate weight $w=-1$) and conductor $n>1$, 
from  \eqref{prop12.2} we have 
\begin{equation}\label{13.5eq1}
\begin{split}
\mathscr{L}_{\I,\boldsymbol{\xi}}^{\mathrm{alg}}(\nu,\hat\phi^{-1})
&=
\mathrm{sp}_{\nu,\hat\phi}\left(\mathcal{L}_{\omega_\I}^{\widetilde{\Gamma}_\infty}(\res_\mathfrak{P}(\mathfrak{Z}_{\boldsymbol{\xi}}))\right)\\
&=
\mathcal{E}(\hat\phi^{-1},\nu)
\cdot
(\omega_{\mathcal{F}_\nu}\otimes\phi^{-1})\left(\log(\mathrm{sp}_{\nu,\hat\phi}(\res_{\mathfrak{P}}(\mathfrak{Z}_{\boldsymbol{\xi}}))\right).\end{split}
\end{equation} 
Using that the characters $\xi_\nu$ and $\phi$ has conductors $p^m$ and $p^n$ respectively, and $n\geq m$, by  \eqref{eq13.2} we have
{\footnotesize{\begin{equation}\label{13.5eq2}
\begin{split}
\mathscr{L}_{\I,\boldsymbol{\xi}}^{\mathrm{alg}}(\nu,\hat\phi^{-1})
&=
\mathcal{E}(\hat\phi^{-1},\nu)\sum_{\sigma\in\Gal(H_{cp^n}/H_c)}(\hat\xi_\nu^{-1}\hat\phi^{-1})(\sigma)
\log(\mathrm{sp}_{\nu}(\res_{\mathfrak{p}}(\cor_{H_c/K}(\mathfrak{X}_{cp^\infty}^\sigma))))(\omega_{\mathcal{F}_\nu^*})
\\
&=
\mathcal{E}(\hat\phi^{-1},\nu)\sum_{\sigma\in\Gal(H_{cp^n}/K)}(\hat\xi_\nu^{-1}\hat\phi^{-1})(\sigma)
\log(\mathrm{sp}_{\nu}(\res_{\mathfrak{p}}(\mathfrak{X}_{cp^\infty}^\sigma)))(\omega_{\mathcal{F}_\nu^*})
\\
&=
\mathcal{E}(\hat\phi^{-1},\nu)\sum_{\sigma\in\Gal(H_{cp^n}/K)}\nu(\mathbf{a}_p)^{-n}(\hat\xi_\nu^{-1}\hat\phi^{-1})(\sigma)
\log(\mathrm{sp}_{\nu}(\res_{\mathfrak{p}}(\mathfrak{X}_{cp^n}^\sigma)))(\omega_{\mathcal{F}_\nu^*})
\\
&=
\mathcal{E}(\hat\phi^{-1},\nu)\left(\frac{p}{\nu(\mathbf{a}_p)}\right)^m\sum_{\sigma\in\Gal(H_{cp^n}/K)}
\nu(\mathbf{a}_p)^{-n}(\hat\xi_\nu^{-1}\hat\phi^{-1})(\sigma)\log(\varrho_m(Q_{cp^{n},m}^\sigma))(\omega_{\mathcal{F}_\nu^*})
\\
&=
\mathcal{E}(\hat\phi^{-1},\nu)\left(\frac{p}{\nu(\mathbf{a}_p)}\right)^m\sum_{\sigma\in\Gal(H_{cp^n}/K)}\nu(\mathbf{a}_p)^{-n}
(\hat\xi_\nu^{-1}\hat\chi_\nu\hat\phi^{-1})(\sigma)\log(\varrho_m( P_{cp^{n+m},m}^\sigma))(\omega_{\mathcal{F}_\nu^*}).
\end{split}\end{equation}}}
Let $F_\infty^*$ be the Coleman primitive of $\omega_{\mathcal{F}_\nu^*}$ on $\mathcal{W}_\infty(p^m)$ which vanishes at $x_\infty$. 
It follows from \eqref{eq9.14} that 
\[\log(\varrho_m( P_{cp^{n+m},m}^\sigma))=\log(j_m( P_{cp^{n+m},m}^\sigma)).\]
Applying Corollary \ref{CoroCole} (and using linearity) we thus obtain 
\[\mathscr{L}_{\I,\boldsymbol{\xi}}^{\mathrm{alg}}(\nu,\hat\phi^{-1})=\mathcal{E}(\hat\phi^{-1},\nu)\left(\frac{p}{\nu(\mathbf{a}_p)}\right)^m\sum_{\sigma\in\Gal(H_{cp^n}/K)}\nu(\mathbf{a}_p)^{-n}
(\hat\xi_\nu^{-1}\hat\chi_\nu\hat\phi^{-1})(\sigma)F_\infty^*({P}_{cp^{n+m},m}^\sigma).\]
On the other hand, since ${P}_{cp^{n+m},m}$ is defined over the subfield $H_{cp^{n-1}}(\zeta_{p^n})$ of $L_{cp^n}$, and $\chi_\nu$ is a primitive character modulo $p^n$,
we see that, after setting $\varphi=\nu(\mathbf{a}_p)^{-n}\hat\xi_\nu^{-1}\hat\chi_\nu\hat\phi^{-1}$ to simplify the notation,  
\begin{equation} \begin{split}
\sum_{\sigma}
\varphi(\sigma)F_\infty^*({P}_{cp^{n+m},m}^\sigma)&= 
\sum_{\sigma}
\varphi(\sigma)F_\infty^*({P}_{cp^{n+m},m}^\sigma)-\frac{\nu(\mathbf{a}_p)}{p}\sum_\sigma\varphi(\sigma)F^*_\infty(\phi({P}_{cp^{n+m},m}^\sigma))\\
&=\sum_{\sigma}
\varphi(\sigma)\Pi(\phi^*)F_\infty^*({P}_{cp^{n+m},m}^\sigma)\\
&=\sum_{\sigma}
\varphi (\sigma) d^{-1}{\mathcal{F}_{\nu , x_\infty}^{[p]}}({P}_{cp^{n+m},m}^\sigma)\\
\end{split}\end{equation}
where the sum is over all $\sigma\in\Gal(H_{cp^n}/K)$, and the last equation follows from \eqref{Col-dminus1} and the fact that $d^{-1}\omega_{\mathcal{F}^{[p]}}=d^{-1}\omega_{{\mathcal{F}^*}^{[p]}}$. Therefore, 
{\footnotesize{
\begin{equation}\label{13.5eq4}\mathscr{L}_{\I,\boldsymbol{\xi}}^{\mathrm{alg}}(\nu,\hat\phi^{-1})=\mathcal{E}(\hat\phi^{-1},\nu)\left(\frac{p}{\nu(\mathbf{a}_p)}\right)^m\cdot\sum_{\sigma\in\Gal(H_{cp^n}/K)}\nu(\mathbf{a}_p)^{-n}(\hat\xi_\nu^{-1}\hat\chi_\nu\hat\phi^{-1})(\sigma)d^{-1}{\mathcal{F}_{\nu, x_\infty}^{[p]}}({P}_{cp^{n+m},m}^\sigma).\end{equation} }}
We now observe that 
\begin{equation}\label{13.5eq3} U_pF_\infty^*=\left(\frac{\nu(\mathbf{a}_p)}{p}\right)F_\infty^*.\end{equation}
Since ${P}_{cp^{n+m},m}=U_p^m{P}_{cp^n,m}=U_p^m x_{cp^n,m}(1)$, it follows from \eqref{13.5eq3} and \eqref{13.5eq4} that 
(use Shimura's reciprocity law to keep trace of the Galois action) 
\begin{equation}\label{13.5eq5}\mathscr{L}_{\I,\boldsymbol{\xi}}^{\mathrm{alg}}(\nu,\hat\phi^{-1})=\mathcal{E}(\hat\phi^{-1},\nu)\sum_{\sigma\in\Gal(H_{cp^n}/K)}\nu(\mathbf{a}_p)^{-n}(\hat\xi_\nu^{-1}\hat\chi_\nu\hat\phi^{-1})(\sigma)d^{-1}{\mathcal{F}_{\nu , x_\infty}^{[p]}}(x_{cp^{n},m}(1)^\sigma).\end{equation}
Since $\Phi_\nu=\nu(\mathbf{a}_p)({\xi}_{\nu, \mathfrak{p}}(p)p)^{-1}$, we have 
\[
\mathcal{E}(\hat\phi^{-1},\nu)
=\frac{\epsilon(\phi)\nu(\mathbf{a}_p)^n}
{{\xi}_{\nu,\mathfrak{p}}(p^n)\cdot p^n}
\] and the result follows. 
\end{proof}

\end{document}
